\section{Introduction}\label{sec:Introduction}

\paragraph{Motivation.}
Many scientific data sets are sequences of functions: examples include trajectories of partial differential equations, spatial fields observed over time, demand curves, spectra, and continuously sampled signals. Treating each observation as a function preserves geometry and regularity that may be obscured by flattening a discretisation into a vector. Functional data analysis provides the statistical language for this viewpoint \cite{RamsaySilverman2005}, while functional autoregressive models and functional time-series methods provide probabilistic foundations for prediction under temporal dependence \cite{Bosq2000,HormannKokoszka2010,AueEtAl2015}. Alongside this classical theory, transformers provide an expressive finite-context mechanism for processing several recent observations. The present thesis studies how these perspectives can be combined and how transformers can be defined so that each functional context determines a conditional probability law for the next function and recursive sampling generates a stochastic process. This raises questions that are not present in a finite-dimensional implementation: the conditional output law must be mathematically admissible and measurable, the shift-and-append update must define a well-posed Markov chain, and its invariant behaviour and moments must be controlled. A separate statistical question is how the transition risk of a fixed model parameter can be estimated from dependent functional observations.

\paragraph{Literature review.}
Probabilistic sequence models describe sequential data through conditional distributions for future observations given their preceding context. Autoregressive transformers use self-attention to construct context-dependent representations of the preceding sequence \cite{VaswaniEtAl2017}. At scale, they exhibit in-context learning, where predictions adapt to information or examples supplied in the context without updating the model parameters \cite{BrownEtAl2020}. Theoretical accounts include Bayesian interpretations \cite{XieEtAl2022}, learnable function classes \cite{GargEtAl2022,AkyurekEtAl2023}, gradient-based and algorithmic interpretations \cite{VonOswaldEtAl2023,LiEtAl2023ICL,ChengEtAl2024}, and information-theoretic decompositions \cite{JeonEtAl2024}. Markovian analyses relate attention to context-conditioned Markov models \cite{IldizEtAl2024}, study induction heads on Markov-chain prompts \cite{EdelmanEtAl2024}, and, most closely, represent finite-context language generation as a finite-state Markov chain by taking the full token context as the state \cite{ZekriEtAl2025}. This state augmentation also underlies the present thesis, but the function-valued setting requires Hilbert-space admissibility and measurability arguments that do not arise in the same form in the finite-state construction.

In parallel, functional data analysis and functional time-series theory study function-valued observations under temporal dependence \cite{RamsaySilverman2005,Bosq2000,HormannKokoszka2010,AueEtAl2015}. Operator learning approximates maps between function spaces; for example, DeepONet, the Fourier neural operator, and neural-operator theory develop this perspective \cite{LuEtAl2021,LiEtAl2021FNO,KovachkiEtAl2023}. The mathematical, numerical foundations, and statistical theory have also been investigated \cite{BoulleTownsend2024,ReinhardtEtAl2026}. More recent developments have sought to combine operator learning with attention and in-context learning, enabled in part by continuum formulations of attention \cite{CalvelloEtAl2025}. Attention-based variants include Galerkin and neural operator transformers \cite{Cao2021,HaoEtAl2023}, while in-context operator learning uses functional demonstrations to infer operators without parameter updates, including continuum-transformer interpretations as operator gradient descent \cite{YangEtAl2023ICON,YangOsher2024,MishraEtAl2026}. However, probabilistic sequence modelling in function space raises a number of open questions concerning well-posedness, recursive generation, and long-term stochastic behaviour.

A separate literature develops concentration tools for dependent observations, including Marton-coupling bounded-difference inequalities \cite{Paulin2015} and transportation inequalities for dependent sequences \cite{DjelloutGuillinWu2004}. These provide general probabilistic machinery rather than model-specific generalisation guarantees. In the present setting, this leaves the statistical question of how dependence in the pretraining data affects transition-risk estimation, separately from the long-term behaviour of the recursively generated process.

\paragraph{Research questions.}
The thesis addresses four questions.
\begin{enumerate}[label=(Q\arabic*)]
    \item How can a functional autoregressive transformer be defined so that every context produces a valid conditional probability measure and the resulting shift-and-append update is a Borel Markov transition kernel?

    \item Which structural properties of the architecture are sufficient to establish existence and uniqueness of an invariant distribution, uniform geometric convergence in total variation, and finite invariant moments for recursive generation?

    \item For a fixed model parameter, how can the gap between empirical and population pretraining risk be controlled when the functional observations are dependent, and what explicit sample-complexity guarantees follow under zero empirical risk?

    \item How do coupling-based and metric-transport approaches differ in their assumptions and applicability across independent trajectories, overlapping windows, mixing chains, and continuous function-valued dynamics?
\end{enumerate}

\paragraph{Main contributions.}
The principal contributions are as follows.
\begin{enumerate}
    \item A functional autoregressive transformer is defined on a Hilbert space, with output heads lying in a fixed finite-dimensional approximation space, that produce basis coefficients specifying the conditional mean and covariance of the Gaussian next-function distribution.

    \item The conditional Gaussian measure is proved to exist uniquely for every context, to have common support, and to depend measurably on the context. Consequently, recursive generation is formulated rigorously as a time-homogeneous Markov chain on the full functional context space.

    \item Uniform conditional moments yield a geometric Foster--Lyapunov drift inequality, while the common approximation space and covariance floor yield a uniform lower bound for the output densities. After the original context has been shifted out, these bounds give a global multistep minorisation. A direct Doeblin argument then proves a unique invariant distribution, uniform geometric convergence in total variation from every initial distribution, support on the product approximation space, and a finite weighted second moment.

    \item For a fixed parameter value, a squared Wasserstein loss between reference and model transition kernels is introduced and shown to reduce exactly to the discrepancy between their conditional output laws. Under a uniform second-moment assumption, Marton coupling and dependent bounded differences give a high-probability generalisation bound and an explicit pretraining sample-complexity bound under zero empirical risk.

    \item A complementary generalisation theory is developed for metric-dependent functional trajectories. A path-space transportation inequality and Wasserstein stability of the reference and model kernels yield concentration for Lipschitz empirical losses and a second zero-risk sample-complexity bound. The architectural Lipschitz estimates required for this result are proved in the appendix, and the thesis distinguishes these pretraining assumptions from the ergodicity of the recursively generated inference chain.
\end{enumerate}

\paragraph{Organisation of the thesis.}
Section~\ref{sec:Preliminaries} introduces the mathematical tools used throughout and defines the functional autoregressive transformer. Section~\ref{sec:General state-space Markov chain framework for functional autoregressive transformers} establishes the Markov-chain formulation, Gaussian-kernel well-posedness, and the Chain's geometric ergodicity. Section~\ref{sec:Fixed-parameter pretraining generalisation} develops the transition loss, the Marton-coupling generalisation and sample-complexity results, the transportation-based metric theory, and the applicability of both approaches to different pretraining-data regimes. Section~\ref{sec:Conclusion and future work} summarises the results and outlines directions for extending the model and theory. Appendix~\ref{app:Lipschitz regularity} proves the global Lipschitz estimates used in the metric generalisation analysis.